\DocumentMetadata{}
\documentclass[12pt]{article}
\usepackage[margin=1in]{geometry}
\usepackage[utf8]{inputenc}
\usepackage[T1]{fontenc}
\usepackage{enumitem}
\usepackage{hyperref}
\usepackage{xurl}
\usepackage{parskip}
\setlength{\parskip}{6pt plus 2pt minus 1pt}
\usepackage{titlesec}
\usepackage{xcolor}
\usepackage{fancyhdr}
\usepackage{tcolorbox}
\tcbuselibrary{breakable}

% UofSC Color Definitions
\definecolor{garnet}{RGB}{115,0,10}
\definecolor{rose}{RGB}{204,46,64}
\definecolor{atlantic}{RGB}{70,106,159}
\definecolor{congaree}{RGB}{31,65,77}
\definecolor{horseshoe}{RGB}{101,120,11}
\definecolor{neutral90}{RGB}{54,54,54}

% Page setup
\pagestyle{fancy}
\fancyhf{}
\lhead{Daily Gamecock \textit{Opinion}}
\rhead{Opening Hooks}
\cfoot{\thepage}
\renewcommand{\headrulewidth}{0.4pt}

% Section formatting
\titleformat{\section}{\Large\bfseries\color{garnet}}{\thesection}{1em}{}
\titleformat{\subsection}{\large\bfseries\color{neutral90}}{\thesubsection}{1em}{}

% Hyperlink setup
\hypersetup{
    colorlinks=true,
    urlcolor=atlantic,
    linkcolor=atlantic,
    citecolor=atlantic
}

% Define tcolorbox environments
\newtcolorbox{tipbox}[1][]{
    colback=congaree!10,
    colframe=congaree,
    boxrule=1pt,
    sharp corners,
    left=8pt,
    right=8pt,
    top=8pt,
    bottom=8pt,
    title={\textbf{#1}},
    fonttitle=\bfseries,
    coltitle=white,
    colbacktitle=congaree,
    toptitle=4pt,
    bottomtitle=4pt,
    breakable,
    parbox=false
}

\newtcolorbox{promptbox}{
    colback=atlantic!10,
    colframe=atlantic,
    boxrule=1pt,
    sharp corners,
    left=8pt,
    right=8pt,
    top=8pt,
    bottom=8pt,
    breakable,
    parbox=false
}

\newtcolorbox{seasonbox}[1][]{
    colback=horseshoe!10,
    colframe=horseshoe,
    boxrule=1pt,
    sharp corners,
    left=8pt,
    right=8pt,
    top=8pt,
    bottom=8pt,
    title={\textbf{#1}},
    fonttitle=\bfseries,
    coltitle=white,
    colbacktitle=horseshoe,
    toptitle=4pt,
    bottomtitle=4pt,
    breakable,
    parbox=false
}

\newtcolorbox{keybox}{
    colback=garnet!5,
    colframe=garnet,
    boxrule=1pt,
    sharp corners,
    left=8pt,
    right=8pt,
    top=8pt,
    bottom=8pt,
    breakable,
    parbox=false
}

% Example counter
\newcounter{examplenum}
\newcommand{\totalexamples}{10}
% #1 = title, #2 = author, #3 = date, #4 = url
\newcommand{\exampleheader}[4]{%
    \newpage
    \stepcounter{examplenum}%
    \noindent{\large\bfseries\color{garnet} Example \theexamplenum\ of \totalexamples}\\[6pt]
    \textbf{Title:} #1\\
    \textbf{Author:} #2\\
    \textbf{Date:} #3\\
    \textbf{Link:} \url{#4}\\[4pt]
    \noindent\rule{\textwidth}{0.4pt}
    \vspace{6pt}
}

% Custom compact title
\newcommand{\doctitle}{%
    \noindent{\LARGE\bfseries Daily Gamecock \textit{Opinion}}\\[2pt]
    {\LARGE\bfseries Opening Hooks}\\[4pt]
    {\normalsize J.C. Vaught. Spring 2026}\\[-2pt]
    \noindent\rule{\textwidth}{1pt}
    \vspace{8pt}
}

\begin{document}
\thispagestyle{plain}

\doctitle

\begin{keybox}
A hook is the first 1--3 sentences of an opinion piece. Its job is to make someone who doesn't care about your topic start reading anyway. This guide covers seven hook types, shows what works and what doesn't using real Daily Gamecock pieces, and teaches you how to match your hook to your argument.
\end{keybox}

\section*{What Is a Hook?}

The first 1--3 sentences of an opinion piece. Its job is to make someone who doesn't care about your topic start reading anyway.

A hook is not a thesis. It's an invitation.

\section*{Why Hooks Matter in Opinion Writing}

\begin{keybox}
News writing has built-in hooks --- the event itself. A fire, an election, a budget cut. The reader already has a reason to care.

Opinion writing doesn't have that luxury. Your reader already knows you have an opinion. Everyone has opinions. The hook has to make them curious about \textit{your} take --- why your argument is worth three minutes of their attention.

Without a strong hook, readers bounce before they ever reach your nutgraf. The piece might be brilliant from paragraph two onward, but nobody will know.
\end{keybox}

\begin{tipbox}[Remember]
If your nutgraf is the reason your piece exists, the hook is the reason anyone reads long enough to find it.
\end{tipbox}

\section*{The Seven Hook Types}

\begin{tipbox}[The Seven Hook Types]
\begin{enumerate}[leftmargin=*]
    \item \textbf{Anecdote} --- A short, specific story about a real person or moment that draws the reader into the world of the piece.
    \item \textbf{Scene-Setting} --- Sensory details or fragments that paint a vivid picture before the argument arrives.
    \item \textbf{Statistic/Data} --- A surprising number or fact that reframes the reader's assumptions.
    \item \textbf{Provocation} --- A bold, assertive claim designed to make the reader react before they think.
    \item \textbf{Rhetorical Question} --- A question posed to the reader that has no obvious answer, creating curiosity.
    \item \textbf{Direct Address} --- Speaking directly to a specific audience, creating intimacy and urgency.
    \item \textbf{News Peg} --- Anchoring the opening to a recent event or decision that gives the piece timeliness.
\end{enumerate}
\end{tipbox}

\section*{Matching Hook to Piece}

Not every hook works for every piece. Here's a rough guide for pairing hook types with common opinion formats:

\begin{itemize}[leftmargin=*]
    \item \textbf{Personal columns} --- Anecdote or scene-setting. You're telling a story rooted in experience; let the reader into that experience from the first sentence.
    \item \textbf{Policy arguments} --- Statistic/data or news peg. You're making a case about systems; ground it in evidence or timeliness right away.
    \item \textbf{Cultural commentary} --- Provocation or rhetorical question. You're challenging how people think; start by disrupting their assumptions.
    \item \textbf{Editorials} --- News peg or direct address. You're speaking on behalf of a position with institutional weight; anchor it in the moment or speak directly to the people affected.
\end{itemize}

These are starting points, not rules. A policy argument can open with an anecdote. A personal column can open with a statistic. The question is always the same: does this hook make someone who doesn't care start reading?

\vspace{12pt}
\begin{center}
    {\Large\bfseries\color{garnet} Examples from the Archive}
\end{center}

% ============================================================
% EXAMPLE 1
% ============================================================
\exampleheader
{Column: Parking and the illusion of equal access}
{JC Vaught}
{4/22/25}
{https://www.dailygamecock.com/article/2025/04/column-parking-and-the-illusion-of-equal-access-opinion-vaught}

\subsection*{Hook Type: Anecdote (Strong)}

\begin{promptbox}
``Rodrigo Rojas, a fourth-year mechanical engineering student, starts every weekday before dawn by plotting what he calls `the loop' --- a three-street circuit that snakes past Assembly, Main and Bull streets in hopes of nabbing one of Columbia's curbside spots. If no spaces appear, his 1999 Mazda Miata, bought for its low cost and 30-mpg efficiency, ends up tucked illegally behind an apartment complex, where he prays the meter maid will miss him until his classes end.''
\end{promptbox}

\textbf{Analysis:} This is a textbook anecdote hook. A specific person, a specific routine, specific details (1999 Miata, 30 mpg, ``the loop''). The reader doesn't know what the column is about yet, but they're already invested in whether Rodrigo gets a parking spot.

The details do double duty --- they establish character AND quietly set up the policy argument (low-cost car = budget-conscious student).

Strong example, no revision needed. Note how the anecdote is tight --- two sentences, not five.

% ============================================================
% EXAMPLE 2
% ============================================================
\exampleheader
{Column: Social media is a culprit of holiday loneliness}
{Christina Ceniccola}
{11/21/21}
{https://www.dailygamecock.com/article/2021/11/column-social-media-is-a-culprit-of-holiday-loneliness-opinion-ceniccola}

\subsection*{Hook Type: Scene-Setting (Strong)}

\begin{promptbox}
``Neighborhood drives through an endless sea of Christmas lights. Grandma's classic pumpkin pie. Snuggling your pet by the fireplace, hot cocoa in hand. It's the most wonderful time of the year \ldots\ right?''
\end{promptbox}

\textbf{Analysis:} This uses sensory fragments to paint a cozy holiday scene, then pulls the rug with ``right?'' The fragments are a stylistic choice --- they feel like scrolling through an Instagram feed, which foreshadows the social media argument.

The hook works because it builds comfort before questioning it. Strong example. The only risk with fragment hooks is overuse --- use them when the rhythm serves the argument, not as a default.

% ============================================================
% EXAMPLE 3
% ============================================================
\exampleheader
{Column: It's time to take politics out of organization funding}
{JC Vaught}
{4/16/25}
{https://www.dailygamecock.com/article/2025/04/column-its-time-to-take-politics-out-of-organization-funding-opinion-vaught}

\subsection*{Hook Type: Personal Anecdote (Strong)}

\begin{promptbox}
``Two years ago, I walked into the Russell House senate chamber as a newly elected senator, later serving as treasurer, believing we could tighten a few bolts and fix what was broken. I was wrong.''
\end{promptbox}

\textbf{Analysis:} This is a first-person variant of the anecdote hook. It works because of the two-word paragraph: ``I was wrong.'' That sentence creates immediate tension --- the reader needs to know what was wrong and why.

The setup (elected senator, optimistic) and the reversal (I was wrong) take only two sentences. Personal anecdotes are powerful in opinion writing because they establish credibility --- the writer isn't just opining from the outside, they were inside the system. But they only work when the personal experience is directly relevant to the argument.

% ============================================================
% EXAMPLE 4
% ============================================================
\exampleheader
{Column: College football playoff needs to expand}
{Jacob Phillips}
{11/30/21}
{https://www.dailygamecock.com/article/2021/12/column-college-football-playoff-needs-to-expand-opinion-phillips}

\subsection*{Hook Type: News Peg --- attempting statistic/data (Weak)}

\begin{promptbox}
``The College Football Playoff (CFP) management committee is set to consider an expansion to an eight-team, or even better, a 12-team playoff format on Dec. 1. The committee should vote to adopt the 12-team format.''
\end{promptbox}

\textbf{Analysis:} This is a common mistake --- the writer treats a news event as a hook, but there's nothing hooky about it. It reads like a news brief, not an opinion opening. The ``even better'' parenthetical reveals the writer's stance too early and too casually.

A data hook should surprise the reader with a number they didn't expect. This piece would benefit from leading with a striking stat about how few teams have ever made the 4-team playoff, or how many undefeated teams have been excluded.

\begin{seasonbox}[Option A --- Statistic + Absurdity]
In 28 years of the BCS and CFP eras, only 13 different programs have played for a national championship. An undefeated UCF team went 25--0 across two seasons and never got an invitation. A one-loss Alabama team that didn't win its own division got in twice. If the College Football Playoff were a restaurant, it would seat the same four families every night and call it a reservation system.
\end{seasonbox}

\begin{seasonbox}[Option B --- Statistic + Scene Contrast]
In 28 years of the BCS and CFP eras, only 13 different programs have played for a national championship. On Jan.\ 10, 2022, Georgia and Alabama will play for the title. Again. Somewhere in Cincinnati, a 13--0 Bearcats team will watch from home --- the best team in the country that the system was designed to exclude.
\end{seasonbox}

\begin{seasonbox}[Option C --- Statistic + News Peg]
In 28 years of the BCS and CFP eras, only 13 different programs have played for a national championship. On Dec.\ 1, the CFP management committee will vote on whether to finally change that --- or keep seating the same four families at the table and calling it a meritocracy.
\end{seasonbox}

% ============================================================
% EXAMPLE 5
% ============================================================
\exampleheader
{Opinion: Movie remakes are lazy, artistically bankrupt}
{Joseph Will}
{11/26/18}
{https://www.dailygamecock.com/article/2018/11/movie-remakes-are-lazy-artistically-bankrupt}

\subsection*{Hook Type: Provocation --- attempted (Weak)}

\begin{promptbox}
``When I saw the trailer for the upcoming `Lion King' remake, I couldn't help but wonder what the point is. The original movie won't be improved upon in any meaningful way, and this hollow CGI cash-grab is going to make a billion dollars regardless of its quality.''
\end{promptbox}

\textbf{Analysis:} The instinct is right --- provocation hooks should make the reader react (``wait, really?''). But ``I couldn't help but wonder'' is passive and hedging. A provocation hook should provoke, not wonder. The second sentence is stronger --- ``hollow CGI cash-grab'' has energy.

The problem is that the writer buried the provocation behind a first-person filter. Drop the ``I'' and lead with the punch.

\begin{seasonbox}[Option A --- Blunt Provocation]
The Lion King remake is going to make a billion dollars, and it doesn't deserve a single one of them. Disney isn't reviving a classic --- it's photocopying one, replacing hand-drawn warmth with hollow CGI, and banking on your nostalgia to cover the difference.
\end{seasonbox}

\begin{seasonbox}[Option B --- Rhetorical Question]
Name one scene in the original \textit{Lion King} that needs better graphics. You can't --- because the film's power was never about photorealism. It was about color, exaggeration, and the way Mufasa's ghost filled an entire sky with hand-painted stars. Disney's remake strips all of that out and calls it an upgrade.
\end{seasonbox}

\begin{seasonbox}[Option C --- Data Punch]
Disney has released 18 live-action remakes of animated classics since 2010. Combined, they've grossed over \$10 billion. Combined, they've contributed zero original ideas. The \textit{Lion King} remake is next in line --- and the most egregious yet, because the original never needed saving.
\end{seasonbox}

% ============================================================
% EXAMPLE 6
% ============================================================
\exampleheader
{Opinion: College lifestyle is detrimental}
{Clara Bergeson}
{12/1/19}
{https://www.dailygamecock.com/article/2019/12/college-lifestyle-is-detrimental}

\subsection*{Hook Type: Rhetorical Question (Mixed)}

\begin{promptbox}
``Why is it that college has become a game of `who has it worse?' From nightly pizza rolls to all-nighters, the idealism of what college health is will forever be one of the most detrimental parts of college culture.''
\end{promptbox}

\textbf{Analysis:} The question is genuinely interesting --- ``who has it worse?'' captures a real campus dynamic. But the second sentence immediately deflates it with vague abstraction (``the idealism of what college health is'').

Rhetorical questions work as hooks when the answer isn't obvious and the next sentence sharpens the question rather than abandoning it. Here, the writer should stay in the question longer.

\begin{seasonbox}[Option A --- Anaphora Buildup]
Why is it that college has become a game of ``who has it worse?'' We brag about all-nighters like they're achievements. We eat pizza rolls for dinner and call it self-care. We treat exhaustion as proof that we're working hard enough. Somewhere along the way, self-destruction became the unofficial curriculum.
\end{seasonbox}

\begin{seasonbox}[Option B --- Scene Flip]
It's 2 a.m.\ in Thomas Cooper Library. A student on her third energy drink high-fives a friend pulling his second all-nighter this week. ``I've slept six hours since Monday,'' she says --- not as a complaint, but as a flex. On this campus, running yourself into the ground is a badge of honor, and nobody's asking why.
\end{seasonbox}

\begin{seasonbox}[Option C --- Statistic Gut-Punch]
Seventy percent of USC students report not getting enough sleep. Forty percent skip meals during finals. And when you tell someone you pulled an all-nighter, the most common response isn't ``are you okay?'' --- it's ``same.'' College hasn't just normalized unhealthy habits. It's turned them into a competition.
\end{seasonbox}

% ============================================================
% EXAMPLE 7
% ============================================================
\exampleheader
{Column: Dear young liberals}
{Linden Atelsek}
{11/30/16}
{https://www.dailygamecock.com/article/2016/11/dear-young-liberals}

\subsection*{Hook Type: Direct Address (Strong)}

\begin{promptbox}
``Dear young liberals,

Many of us are not used to losing. That's probably why waking up on Nov. 9 felt so different from waking up on Nov. 8, although we were living in the exact same country as we were the day before.''
\end{promptbox}

\textbf{Analysis:} The epistolary format immediately signals intimacy and urgency --- this isn't a column, it's a letter. ``Many of us are not used to losing'' is disarmingly honest and avoids the self-righteous tone that plagued post-2016 opinion writing.

The observation that Nov. 9 ``felt different'' from Nov. 8 despite being ``the exact same country'' is sharp --- it names the emotional experience without melodrama. Strong example of direct address done with restraint.

% ============================================================
% EXAMPLE 8
% ============================================================
\exampleheader
{Opinion: Don't fool yourself about class mobility}
{Stephanie Woronko}
{12/9/17}
{https://www.dailygamecock.com/article/2017/12/dont-fool-yourself-about-class-mobility}

\subsection*{Hook Type: Pop Culture Entry (Mixed)}

\begin{promptbox}
``In between catchy songs and talking animals, children are drawn to rags-to-riches stories like Aladdin and Cinderella. From professional athletes to CEOs to The Little Engine That Could, we learn that hard work and a little luck are all you need to climb the social ladder.''
\end{promptbox}

\textbf{Analysis:} The pop culture references (Aladdin, Cinderella, The Little Engine That Could) are accessible and disarming --- the reader expects a light cultural observation before the argument pivots to economic reality. This is effective because the fairy tales ARE the argument: the hook demonstrates the myth the piece will dismantle.

The risk with pop culture hooks is dating your piece --- Aladdin works because it's timeless. A reference to a trending meme won't. Good example, but the second sentence (``hard work and a little luck'') states the thesis too early --- let the fairy tale imagery breathe before pivoting.

% ============================================================
% EXAMPLE 9
% ============================================================
\exampleheader
{Opinion: esports legitimize gaming}
{Rodney Davis}
{12/8/17}
{https://www.dailygamecock.com/article/2017/12/esports-legitimize-gaming}

\subsection*{Hook Type: Sensory/Experiential --- second-person immersion (Strong)}

\begin{promptbox}
``Crowds are roaring next to you, adrenaline is pumping and your throat is aching from how much you have screamed. You cannot imagine being at a more heart-pumping event. Everyone loves the feeling of cheering for their team alongside thousands of strangers.''
\end{promptbox}

\textbf{Analysis:} The writer drops the reader into a physical experience --- roaring, aching throat, adrenaline. The trick is that the reader assumes this is a football game or basketball arena, and the piece is about to reveal it's an esports event. This bait-and-switch only works if the sensory details are vivid enough to create a genuine assumption. Here they are.

The ``you'' perspective forces identification. Strong technique for pieces that challenge the reader's assumptions about what counts as legitimate.

% ============================================================
% EXAMPLE 10
% ============================================================
\exampleheader
{Column: Night football games are better than afternoon kickoffs}
{Sofie Kurzawa}
{4/1/25}
{https://www.dailygamecock.com/article/2025/04/column-night-football-games-are-better-than-afternoon-kickoffs-opinion-kurzawa}

\subsection*{Hook Type: Contrast (Strong)}

\begin{promptbox}
``It's 7 a.m. on gameday in Columbia. You're barely awake before your friend's elaborate pregame plans start to stress you out. And if that's not bad enough, the 90-degree temperatures only make things worse.''
\end{promptbox}

\textbf{Analysis:} This hook works by making the reader feel the misery of the status quo before the argument even arrives. Second person (``you're barely awake'') forces identification. The specific details (7 a.m., 90 degrees, stressed by pregame plans) make it feel lived-in rather than hypothetical.

The implicit contrast --- night games DON'T have these problems --- is already forming in the reader's mind before the writer states it. This is a strong use of contrast: show the problem, let the reader want the solution.

% ============================================================
% BACK MATTER
% ============================================================
\newpage

\section*{Common Hook Mistakes}

\begin{itemize}[leftmargin=*]
    \item Starting with a dictionary definition (``According to Merriam-Webster\ldots'')
    \item Starting with ``In today's society\ldots'' or ``Since the dawn of time\ldots''
    \item Starting with your thesis (that's a nutgraf, not a hook)
    \item Rhetorical questions that have obvious answers (``Don't you think racism is bad?'')
    \item Hooks that have nothing to do with the piece (bait-and-switch without the switch)
    \item Hooks that are longer than 3--4 sentences (you're writing a lede, not a novel)
    \item Starting with ``I'' unless the personal experience is directly relevant
\end{itemize}

\begin{tipbox}[The Hook Test]
Read your first two sentences to someone who doesn't care about your topic. If they say ``okay, and?'' --- your hook works. If they say ``cool'' and look at their phone --- rewrite it.
\end{tipbox}

\newpage

\section*{Practice Exercise}

For each of the following opinion topics, draft \textbf{two different hook types}. Use any of the seven types from this guide, but don't repeat the same type twice for the same topic.

\begin{enumerate}[leftmargin=*]
    \item \textbf{USC dining halls should offer more plant-based options.}

    \vspace{2.5in}

    \item \textbf{Greek life rush should be moved to spring semester.}

    \vspace{2.5in}

    \item \textbf{The university should cancel classes during extreme heat.}

    \vspace{2.5in}
\end{enumerate}

\end{document}
